\comoleerla{cada punto es un $(w_1, w_2)$ de la malla con $w_3 = 1 - w_1 - w_2$; círculo azul, gana Bogotá; triángulo naranja, gana Santander; la línea es la frontera exacta.}
{Bogotá gana en 17 de 66 puntos y la frontera corta $w_3=0$ en $w_1 = 0.4576$.}
{Medido: Bogotá es más corta ($\Delta D = -2.0120$) pero más cara en peaje ($\Delta P = 1.6971$) y en riesgo ($\Delta R = 2.8788$); la frontera es el lugar donde esos tres términos se compensan.}
